%@ pages 387-390
%@ status translated
% Chapter IV, section 2: The structure of injective objects in a locally
% noetherian category (pp. 387--390). Transcribed from the page images
% in reference/source.pdf (PDF pp. 65--68), not from the OCR text layer,
% which mangles subscripts and fraktur letters. The section opens with
% its heading in the middle of p. 387 and runs to the section~3 heading
% in the middle of p. 390 (the end of the proof of Proposition~9 and
% the two closing paragraphs stand at the top of p.~390).
%
% Numbering is automatic (see src/preamble.tex). This section prints
% Proposition 6, Proposition 7, Theorem 2, Proposition 8, one
% corollary printed without a number (hence `unnumberedcorollary')
% and Proposition 9.
%
% Conventions: $\mathbf{KI}_n$ is the homotopy-category construction of
% Chapter~I, \S9 (cf.\ Corollary~1 to Proposition~14); bold capitals
% $\mathbf{L}$, $\mathbf{N}$, $\mathbf{P}$ denote sets (index sets and
% sets of parts).
%
% No emendations.

\section{The Structure of Injective Objects in a Locally Noetherian Category}

--- The statements of the preceding paragraph simplify when
$\mathbf{A}$ is a locally noetherian category. Corollary~1 to
Theorem~1 (Chapter~II) has indeed the following consequence:

\begin{proposition}
Let $\mathbf{A}$ be a locally noetherian category and
$(I_l)_{l\in\mathbf{L}}$ a family of injective objects of
$\mathbf{A}$. The direct sum $\sum_{l\in\mathbf{L}} I_l$ is then an
injective object of $\mathbf{A}$.
\end{proposition}

\begin{proposition}
For every locally noetherian category $\mathbf{A}$, the localizing
subcategory $\mathbf{A}_\omega$ coincides with $\mathbf{A}$; moreover,
the categories
$\mathbf{A}_{\alpha \mathbin{T} 0}/\mathbf{A}_\alpha$ are locally
finite.
\end{proposition}

We know indeed that $\mathbf{A}/\mathbf{A}_\omega$ possesses no simple
object. As $\mathbf{A}$ is a locally noetherian category, the same
holds for $\mathbf{A}/\mathbf{A}_\omega$. If $\mathbf{A}_\omega$ did
not coincide with $\mathbf{A}$, the category
$\mathbf{A}/\mathbf{A}_\omega$ would therefore contain a nonzero
noetherian object $M$. For every maximal proper subobject $N$ of $M$,
the quotient $M/N$ would be simple, which is absurd.

The last assertion follows from the fact that every locally
noetherian category whose Krull dimension is zero is locally finite.

Modulo Propositions~6 and~7, Theorem~1 and Proposition~5 read as
follows:

\begin{theorem}[Matlis~\cite{matlis}]
Let $\mathbf{A}$ be a locally noetherian category. Every injective
$I$ is isomorphic to the direct sum $\sum_{l\in\mathbf{L}} I_l$ of a
family $(I_l)_{l\in\mathbf{L}}$ of indecomposable injective objects.

If $(J_n)_{n\in\mathbf{N}}$ is a second family of indecomposable
injectives, and if the direct sums $\sum_{l\in\mathbf{L}} I_l$ and
$\sum_{n\in\mathbf{N}} J_n$ are isomorphic, there exists a bijection
$h$ from $\mathbf{L}$ onto $\mathbf{N}$ such that $I_l$ is isomorphic
to $J_{h(l)}$.
\end{theorem}

The second assertion of Theorem~2 also follows from Theorem~1
(Chapter~I).

\begin{proposition}[exchange theorem]
Let $\mathbf{A}$ be a locally noetherian category and
$(I_l)_{l\in\mathbf{L}}$ a family of indecomposable injectives of
$\mathbf{A}$. If $J$ is a direct factor of
$\sum_{l\in\mathbf{L}} I_l$, there exists a part $\mathbf{N}$ of
$\mathbf{L}$ such that $\sum_{l\in\mathbf{N}} I_l$ is a complement of
$J$ in $\sum_{l\in\mathbf{L}} I_l$.
\end{proposition}

We suppose in the rest of this paragraph that $\mathbf{A}$ is a
locally noetherian category. We say that the category $\mathbf{A}$ is
\emph{connected} if it is not equivalent to the product of two nonzero
categories.

\begin{unnumberedcorollary}
Every locally noetherian category $\mathbf{A}$ is equivalent to the
product of a family $(\mathbf{A}_n)_{n\in E}$ of connected categories.
\end{unnumberedcorollary}

Let indeed $\mathbf{I}$ be the full subcategory of $\mathbf{A}$ whose
objects are the injective objects of $\mathbf{A}$. By Corollary~1 to
Proposition~14 (Chapter~I), it suffices to show that $\mathbf{I}$ is
the product of a family $(\mathbf{I}_n)_{n\in E}$ of additive
categories such that $\mathbf{KI}_n$ is a connected category for every
$n$. For this we agree to say that two indecomposable injectives $I$
and $J$ are equivalent if there exists a finite sequence $(I_l)$,
$1 \le l \le k$, of indecomposable injectives satisfying the following
conditions: $I_1 = I$, $I_k = J$; if $1 \le l \le k-1$, the abelian
groups $\operatorname{Hom}(I_l, I_{l+1})$ and
$\operatorname{Hom}(I_{l+1}, I_l)$ are not both zero. We denote by $E$
the set of equivalence classes thus obtained (Brauer calls such a
class a \emph{block}; to be precise, let us say that Brauer considers
the indecomposable projective modules over an artinian ring $A$; he
says that two indecomposable projectives $P$ and $Q$ belong to the
same block if there exists a sequence $P_1 = P$, $P_2$, \dots,
$P_k = Q$ formed of indecomposable projectives such that the groups
$\operatorname{Hom}_A(P_l, P_{l+1})$ and
$\operatorname{Hom}_A(P_{l+1}, P_l)$ are never both zero; we shall see
in paragraph~3 that Brauer's theory is a particular case of the one
we develop here). If $n$ is such a class, let $\mathbf{I}_n$ be the
full subcategory of $\mathbf{I}$ defined by the direct sums of
indecomposable injectives of $n$. It is clear that the category
$\mathbf{I}$ is equivalent to the product of the categories
$\mathbf{I}_n$; moreover, $\mathbf{I}_n$ is not equivalent to the
product of two nonzero additive categories; whence the assertion.

Let $\mathbf{A}_n$ be the localizing subcategory of $\mathbf{A}$
defined by the objects $M$ such that $\operatorname{Hom}(M, I)$ is
zero for every indecomposable injective not belonging to $n$. It is
clear that the categories $\mathbf{A}_n$ and $\mathbf{KI}_n$ are
equivalent. The category $\mathbf{A}$ is therefore equivalent to the
product of the connected categories $\mathbf{A}_n$. We shall say that
$\mathbf{A}_n$ is a \emph{connected component} of $\mathbf{A}$.

We end this paragraph with a generalisation of the Lasker-Noether
primary decomposition. For every element $s$ of the spectrum of
$\mathbf{A}$, we choose an indecomposable injective $I_s$ of type $s$.
If $M$ is an object of $\mathbf{A}$, we say that $s$ is
\emph{associated} to $M$ if the injective envelope of $M$ contains a
subobject isomorphic to $I_s$. The set of elements $s$ of
$\operatorname{Sp}(\mathbf{A})$ which are associated to $M$ will be
denoted $\operatorname{ass}(M)$. We say that $M$ is \emph{isotypic of
type $s$} if $\operatorname{ass}(M)$ has $s$ as its only element.

\begin{proposition}
Let $M$ be a noetherian object of a locally noetherian category
$\mathbf{A}$. There exists a map $s \mapsto N_s$ from a finite part
$\mathbf{L}$ of $\operatorname{Sp}(\mathbf{A})$ into the set of
subobjects of $M$ satisfying the following conditions: $M/N_s$ is
isotypic of type $s$; the intersection of the $N_s$ is zero; if
$\mathbf{P}$ is a part of $\mathbf{L}$ distinct from $\mathbf{L}$,
the intersection $\bigcap_{s\in\mathbf{P}} N_s$ is not zero. If these
conditions are satisfied, $\mathbf{L}$ is equal to
$\operatorname{ass}(M)$.
\end{proposition}

Let indeed $i\colon M \to I$ be an injective envelope of $M$. We
suppose for simplicity that $M$ is a subobject of $I$ and that $i$ is
the canonical morphism from $M$ into $I$. We suppose also that $I$ is
the direct sum of a family $(I_l)_{l\in\mathbf{L}}$ of indecomposable
injectives. The subobjects $M \cap I_l$ of $M$ are then nonzero for
every $l$ and their sum is direct. As $M$ is noetherian, the set
$\mathbf{L}$ is finite. If $s$ is an element of
$\operatorname{ass}(M)$, denote by $J_s$ the direct sum of the $I_l$
isomorphic to $I_s$. In these conditions, it suffices to choose
$\mathbf{L}$ equal to $\operatorname{ass}(M)$ and $N_s$ equal to
$M \cap \left(\sum_{t \ne s} J_t\right)$.

Conversely, suppose given $\mathbf{L}$ and the map $s \mapsto N_s$.
The canonical morphisms $q_s\colon M \to M/N_s$ define a monomorphism
$q$ from $M$ into the direct sum $N$ of the $M/N_s$. We deduce the
formula
$\operatorname{ass}(M) \subset \operatorname{ass}(N) = \mathbf{L}$.
Let moreover $M_s$ be the intersection of the $N_t$ for $t$ distinct
from $s$. The morphism $q_s$ induces a monomorphism from $M_s$ into
$M/N_s$. As $\operatorname{ass}(M_s)$ is nonempty and
$\operatorname{ass}(M/N_s) = \{s\}$, we have the following formulas:
$s \in \operatorname{ass}(M_s) \subset \operatorname{ass}(M)$. This
proves that $\mathbf{L}$ is contained in $\operatorname{ass}(M)$,
which remained to be proved.

We leave to the reader the task of transcribing into the language of
this article the many propositions of the primary decomposition. We
also leave to him the task of formulating the dual statements of the
preceding statements.
